% Actividades de síntesis de cada unidad (resumen, mapa conceptual y diccionario científico).
% Se incluye en el anexo de Libro_4_ESO_ES.tex. Ordenado por unidades.

%========================================================================================
\section{Unidad 1. El trabajo científico}
%========================================================================================
% (sacado del final de UD1_Trabajo_cientifico.tex)

\begin{enumerate}[label=\numberballred]
    \item \textbf{Resumen.} Elabora un resumen de la unidad respondiendo a estas preguntas:
    \begin{itemize}
        \item ¿Cuáles son las etapas del método científico? ¿Qué diferencia hay entre causalidad y correlación?
        \item ¿Cómo se analizan los resultados de un experimento? ¿Qué diferencia hay entre una ley, una teoría y un modelo?
        \item ¿Qué son las magnitudes? ¿Cuáles son las unidades básicas del SI? ¿Para qué sirve la notación científica?
        \item ¿Qué son las ecuaciones de dimensiones?
        \item ¿Cómo se expresa el resultado de una medida? ¿Cómo se calculan el error absoluto y el relativo?
        \item ¿Qué diferencia hay entre las magnitudes escalares y las vectoriales? ¿Cómo se suman los vectores?
        \item ¿Qué materiales e instrumentos de laboratorio conoces? ¿Qué normas de seguridad hay que seguir?
    \end{itemize}
    \item \textbf{Mapa conceptual.} Construye un esquema con estos conceptos: \emph{método científico, magnitudes básicas y derivadas, Sistema Internacional, notación científica, ecuación de dimensiones, error absoluto y relativo, magnitudes escalares y vectoriales}.
    \item \textbf{Diccionario científico.} Define con tus palabras: \emph{hipótesis, variable dependiente e independiente, ley, teoría, modelo, tabla, gráfica, interpolar, extrapolar, redondeo, cifra significativa, precisión, vector, módulo, dirección, sentido, vector opuesto}.
\end{enumerate}

%========================================================================================
\section{Unidad 2. El átomo y la Tabla Periódica}
%========================================================================================
% (sacado del final de UD_2_el_atomo_tp.tex)

\begin{enumerate}[label=\numberballblue]
    \item \textbf{Resumen.} Elabora un resumen de la unidad respondiendo a estas preguntas:
    \begin{itemize}
        \item ¿Cuáles son las ideas principales de los modelos de Dalton y de Thomson?
        \item ¿Qué partículas y radiaciones emiten las sustancias radiactivas?
        \item ¿Dónde están las partículas subatómicas en los modelos de Rutherford y de Bohr?
        \item ¿Cuándo es un átomo neutro? ¿Cuándo adquiere carga positiva? ¿Y negativa?
        \item ¿Qué significa la notación \ch{^{A}_{Z}X}? ¿Qué es un isótopo?
        \item ¿Cómo se mide la masa de un átomo? ¿Qué es la masa atómica relativa?
        \item ¿Qué son los espectros de emisión? ¿Por qué no podía explicarlos el modelo de Rutherford?
        \item ¿Cómo explica el modelo de Bohr-Sommerfeld los espectros de emisión?
        \item ¿Cómo se escribe la configuración electrónica de un elemento?
        \item ¿Cómo se ordenan los elementos en el sistema periódico?
        \item ¿Cómo varía la configuración electrónica al avanzar en un período y al bajar en un grupo?
    \end{itemize}
    \item \textbf{Mapa conceptual.} Construye un esquema de la unidad que incluya: \emph{modelo planetario de Rutherford, espectros de emisión, modelo de Bohr, sistema periódico, grupos, períodos, configuración electrónica, elementos representativos, de transición y de transición interna, metales y no metales}.
    \item \textbf{Diccionario científico.} Define con tus palabras: \emph{electrón, protón, neutrón, modelo atómico, radiactividad, partículas alfa y beta, rayos gamma, número atómico, número másico, isótopo, nivel y subnivel de energía, espectro, masa atómica relativa}.
\end{enumerate}

%========================================================================================
\section{Unidad 3. Enlace Químico}
%========================================================================================
% (sacado del final de UD3_enlace_quimico.tex)

\begin{enumerate}[label=\numberballred]
    \item \textbf{Resumen.} Elabora un resumen de la unidad respondiendo a estas preguntas:
    \begin{itemize}
        \item ¿Qué son los electrones de valencia? ¿Qué es un diagrama de Lewis?
        \item ¿Qué es un enlace químico? ¿Qué dice la regla del octeto?
        \item ¿Qué tiene de especial la configuración electrónica de los gases nobles?
        \item ¿Cómo se forman las moléculas de \ch{F2}, \ch{O2} y \ch{N2}? ¿Qué propiedades tienen estas sustancias?
        \item ¿Qué sustancias forman cristales covalentes y qué propiedades tienen?
        \item ¿Qué es un dipolo? ¿Qué fuerzas intermoleculares existen y cómo se producen?
        \item ¿Cómo se forman y qué propiedades tienen los compuestos iónicos?
        \item ¿Cómo se mantienen unidos los átomos de un metal?
        \item ¿Cómo se calcula la masa molecular de una sustancia? ¿Y su composición centesimal?
        \item ¿Qué es la cantidad de sustancia y en qué unidad se mide? ¿Qué es la masa molar?
    \end{itemize}
    \item \textbf{Mapa conceptual.} Construye un esquema de la unidad con estos conceptos: \emph{elementos} (átomos aislados, moléculas covalentes, cristales covalentes y cristales metálicos), \emph{compuestos} (moléculas covalentes, cristales iónicos y cristales covalentes), \emph{masa molecular} y \emph{mol}.
    \item \textbf{Diccionario científico.} Define con tus palabras: \emph{diagrama de Lewis, regla del octeto, enlace químico, gas noble, molécula diatómica, enlace sencillo, doble y triple, mar de electrones, red cristalina, diamante, grafito, dipolo, enlace de hidrógeno, fuerzas de Van der Waals, constante de Avogadro}.
\end{enumerate}

%========================================================================================
\section{Unidad 4. La química del carbono}
%========================================================================================
% (sacado del final de UD4_organica.tex)

\begin{enumerate}[label=\numberballred]
    \item \textbf{Resumen.} Elabora un resumen de la unidad respondiendo a estas preguntas:
    \begin{itemize}
        \item ¿Cuántos electrones comparten dos carbonos unidos por un enlace doble? ¿Y por uno triple?
        \item ¿Qué diferencia hay entre una cadena saturada y una insaturada?
        \item ¿Cuál es el grupo funcional de un alcohol? ¿Y el de una amina?
        \item ¿Qué compuestos del carbono son más solubles en agua: los que tienen grupos polares o los que no los tienen?
        \item ¿Con qué terminación se nombran los alcanos, los alquenos y los alquinos? ¿Qué nombre recibe un hidrocarburo con solo enlaces simples?
        \item ¿Cómo se llaman los grupos funcionales de los alcoholes, los aldehídos, las cetonas y los ácidos carboxílicos? ¿Qué terminación tiene cada uno?
        \item ¿En qué se diferencia una amina primaria de una secundaria y de una terciaria?
        \item ¿Qué sustancias del carbono de especial interés conoces?
    \end{itemize}
    \item \textbf{Mapa conceptual.} Construye un esquema de la unidad con estos conceptos: \emph{estructura tetraédrica del carbono, modelos de bolas y varillas, enlaces simples, dobles y triples, formulación, solubilidad, hidrocarburos, compuestos oxigenados, compuestos nitrogenados, compuestos de interés biológico e industrial}.
    \item \textbf{Diccionario científico.} Define con tus palabras: \emph{cadena saturada e insaturada, fórmula desarrollada, fórmula geométrica, grupo funcional, serie homóloga, radical, alcano, alqueno, alquino, alcohol, grupo hidroxilo, aldehído, cetona, grupo carbonilo, grupo carboxilo, éster, grupo amino, monómero, polímero, termoplástico}.
\end{enumerate}

%========================================================================================
\section{Unidad 5. La materia y los sistemas materiales}
%========================================================================================
% (sacado del final de UD5_Sistema_materiales.tex)

\begin{enumerate}[label=\numberballred]
    \item \textbf{Resumen.} Elabora un resumen de la unidad respondiendo a estas preguntas:
    \begin{itemize}
        \item ¿Qué es la materia? ¿Qué propiedades tienen los tres estados de agregación? ¿Cómo son las fuerzas de cohesión en cada uno?
        \item ¿Qué relación hay entre la presión y el volumen de un gas a temperatura constante? ¿Y entre el volumen y la temperatura a presión constante? ¿Y entre la presión y la temperatura a volumen constante?
        \item ¿Qué relación hay entre la presión, el volumen y la temperatura de un gas? ¿Qué es el volumen molar?
        \item ¿Qué es un sistema material? ¿Cómo se separan los componentes de uno heterogéneo?
        \item ¿Cómo se clasifican los sistemas homogéneos? ¿Cómo se expresa la concentración de una disolución?
        \item ¿Cómo se prepara una disolución de concentración conocida? ¿Cómo se separan sus componentes?
    \end{itemize}
    \item \textbf{Mapa conceptual.} Construye un esquema de la unidad con estos conceptos: \emph{materia, estados de agregación, ley de Boyle-Mariotte, ley de Charles, ley de Gay-Lussac, ley combinada, ley general de los gases, sistema material heterogéneo y homogéneo, sustancia pura, disolución, concentración}.
    \item \textbf{Diccionario científico.} Define con tus palabras: \emph{presión, volumen, temperatura absoluta, condiciones estándar, constante universal de los gases, volumen molar, separación magnética, filtración, decantación, disolución, soluto, disolvente, tanto por ciento en masa y en volumen, molaridad, cristalización, extracción, destilación}.
\end{enumerate}

%========================================================================================
\section{Unidad 6. Las reacciones químicas}
%========================================================================================
% (sacado del final de UD6_reacciones_quimicas.tex)

\begin{enumerate}[label=\numberballred]
    \item \textbf{Resumen.} Elabora un resumen de la unidad respondiendo a estas preguntas:
    \begin{itemize}
        \item ¿Qué es una reacción química? ¿Cuándo es eficaz un choque entre reactivos? ¿Qué es la energía de activación?
        \item ¿Qué dice la ley de conservación de la masa? ¿Cómo se ajusta una ecuación química?
        \item ¿Cómo se hace un cálculo estequiométrico? ¿Qué es el rendimiento?
        \item ¿Qué factores afectan a la velocidad de una reacción? ¿Cómo actúa un catalizador?
        \item ¿Cuándo es una reacción exotérmica o endotérmica?
        \item ¿Qué son las reacciones de síntesis, de oxidación-reducción y de combustión? ¿Qué son los procesos electroquímicos?
        \item ¿Cómo se distingue un ácido de una base? ¿Qué es el pH? ¿Qué es una neutralización?
        \item ¿Qué es el efecto invernadero y cuáles son las consecuencias de su incremento?
        \item ¿Qué etapas tiene un proceso químico industrial?
    \end{itemize}
    \item \textbf{Mapa conceptual.} Construye un esquema de la unidad con estos conceptos: \emph{mecanismo de reacción, teoría de las colisiones, ley de conservación de la masa, ajuste de ecuaciones, estequiometría, reacciones exotérmicas y endotérmicas, velocidad de reacción, síntesis, neutralización, redox, combustión, procesos electroquímicos}.
    \item \textbf{Diccionario científico.} Define con tus palabras: \emph{reactivos, productos, complejo activado, energía de activación, ecuación química, calor de reacción, reacción reversible, catalizador, ácido, base, pH, indicador, electrólisis, pila alcalina, galvanoplastia, metalurgia, efecto invernadero, calentamiento global, cambio climático}.
\end{enumerate}

%========================================================================================
\section{Unidad 7. Los movimientos rectilíneos}
%========================================================================================
% (sacado del final de UD7_movimientos_rectilineos.tex)

\begin{enumerate}[label=\numberballred]
    \item \textbf{Resumen.} Elabora un resumen de la unidad respondiendo a estas preguntas:
    \begin{itemize}
        \item ¿Cuándo decimos que un cuerpo se mueve? ¿Qué es un sistema de referencia? ¿Qué es el vector de posición?
        \item ¿Qué información dan las gráficas posición-tiempo y velocidad-tiempo?
        \item ¿Qué diferencia hay entre trayectoria, espacio recorrido y desplazamiento?
        \item ¿Qué diferencia hay entre velocidad media e instantánea? ¿Qué relación hay entre aceleración y variación de la velocidad?
        \item ¿Qué características tiene un MRU y qué ecuación lo describe? ¿Y un MRUV?
        \item ¿Cómo son las gráficas \(s\)-\(t\), \(v\)-\(t\) y \(a\)-\(t\) de un MRUA y de un MRUR?
        \item ¿En qué se diferencian la caída libre, el lanzamiento hacia abajo y el lanzamiento hacia arriba?
    \end{itemize}
    \item \textbf{Mapa conceptual.} Construye un esquema con estos conceptos: \emph{móvil, posición, sistema de referencia, gráfica posición-tiempo, velocidad media e instantánea, gráfica velocidad-tiempo, aceleración media e instantánea, MRU, MRUV, caída libre, lanzamiento vertical, aceleración de la gravedad}.
    \item \textbf{Diccionario científico.} Define con tus palabras: \emph{vector de posición, trayectoria, desplazamiento, espacio recorrido, vector velocidad, vector aceleración, rapidez media, ecuación de la posición, ecuación de la velocidad, distancia de seguridad, tiempo de reacción}.
\end{enumerate}

%========================================================================================
\section{Unidad 8. Las fuerzas y los cambios de movimiento}
%========================================================================================
% (sacado del final de UD8_fuerzas.tex)

\begin{enumerate}[label=\numberballred]
    \item \textbf{Resumen.} Elabora un resumen de la unidad respondiendo a estas preguntas:
    \begin{itemize}
        \item ¿Qué efectos pueden producir las fuerzas? ¿Cómo se representa una fuerza?
        \item ¿Cómo se componen dos fuerzas de la misma dirección y sentido? ¿Y de sentidos opuestos? ¿Y perpendiculares?
        \item ¿Qué dicen la primera, la segunda y la tercera ley de Newton? ¿Qué es la inercia?
        \item ¿Por qué no se anulan entre sí las fuerzas de acción y reacción?
        \item ¿Qué relación hay entre la masa y el peso de un cuerpo? ¿Dónde está aplicado el peso?
        \item ¿Qué fuerzas actúan sobre un cuerpo que se desliza por un plano inclinado bajo la acción de su propio peso?
        \item ¿Qué características tiene la fuerza de rozamiento y de qué depende? ¿Cómo se dibuja?
    \end{itemize}
    \item \textbf{Mapa conceptual.} Construye un esquema de la unidad con estos conceptos: \emph{representación de fuerzas, composición de fuerzas, primera, segunda y tercera ley de Newton, peso, plano inclinado, fuerza de rozamiento, fuerzas sobre un cuerpo colgado de un hilo, sobre una superficie horizontal y sobre un plano inclinado}.
    \item \textbf{Diccionario científico.} Define con tus palabras: \emph{fuerza, newton, dinamómetro, fuerza resultante, equilibrio, fuerza neta, inercia, acción-reacción, peso, centro de gravedad, fuerza normal, tensión, coeficiente de rozamiento}.
\end{enumerate}

%========================================================================================
\section{Unidad 9. El movimiento circular. La gravedad y otras fuerzas}
%========================================================================================
% (sacado del final de UD9_movimiento_circular.tex)

\begin{enumerate}[label=\numberballred]
    \item \textbf{Resumen.} Elabora un resumen de la unidad respondiendo a estas preguntas:
    \begin{itemize}
        \item ¿Cuáles son las características del MCU? ¿Qué magnitudes angulares conoces y qué relación tienen con las lineales?
        \item ¿Qué relación hay entre el período y la frecuencia?
        \item ¿Por qué un MCU tiene aceleración? ¿Qué ocurre cuando deja de actuar la fuerza centrípeta?
        \item ¿Qué dicen las tres leyes de Kepler?
        \item ¿Cuál es el enunciado de la ley de gravitación universal? ¿Qué características tiene la fuerza gravitatoria?
        \item ¿Cómo explica Newton la caída y el peso de los cuerpos con la ley de gravitación universal?
        \item ¿Qué fuerza mantiene en órbita a la Luna y a los satélites artificiales? ¿Cuál es el período de un satélite geoestacionario?
        \item ¿Cómo calcularías tu peso en la superficie de cualquier astro del sistema solar?
        \item ¿En qué se parecen y en qué se diferencian la ley de Coulomb y la de gravitación universal?
        \item ¿Cómo se manifiestan las fuerzas eléctrica y magnética en la naturaleza? ¿Qué demostraron Oersted y Faraday?
    \end{itemize}
    \item \textbf{Mapa conceptual.} Construye un esquema de la unidad con estos conceptos: \emph{magnitudes angulares y lineales, ecuación del MCU, período y frecuencia, aceleración centrípeta, fuerza centrípeta, leyes de Kepler, ley de gravitación universal, síntesis newtoniana, peso, movimiento de la Luna y de los satélites, ley de Coulomb, potencial eléctrico, imanes, electromagnetismo}.
    \item \textbf{Diccionario científico.} Define con tus palabras: \emph{radio vector, desplazamiento angular, radián, velocidad angular, período, frecuencia, aceleración centrípeta, fuerza centrípeta, perihelio, afelio, constante de gravitación universal, satélite geoestacionario, carga eléctrica, potencial eléctrico, diferencia de potencial, imán permanente, línea de campo magnético, corriente inducida, aurora polar}.
\end{enumerate}

%========================================================================================
\section{Unidad 10. Fuerzas en los fluidos}
%========================================================================================
% (sacado del final de UD10_fluidos.tex)

\begin{enumerate}[label=\numberballred]
    \item \textbf{Resumen.} Elabora un resumen de la unidad respondiendo a estas preguntas:
    \begin{itemize}
        \item ¿Qué es la presión y en qué unidades se mide?
        \item ¿Qué es la presión hidrostática? ¿De qué factores depende?
        \item ¿Qué dice el principio fundamental de la hidrostática? ¿Cómo se calcula la diferencia de presión entre dos puntos de un líquido?
        \item ¿Qué dice el principio de Pascal? ¿Qué relación hay entre las fuerzas y las superficies de los émbolos de una prensa hidráulica?
        \item ¿Qué características tiene la presión atmosférica? ¿Cómo se mide?
        \item ¿Qué es un mapa meteorológico y qué información aporta? ¿Qué relación hay entre las borrascas, los anticiclones y el tiempo?
        \item ¿En qué se diferencian un frente frío y uno cálido?
        \item ¿Qué dice el principio de Arquímedes? ¿Qué relación hay entre el empuje, el peso de un cuerpo y su peso aparente en un líquido?
        \item ¿Cuándo flota un cuerpo y cuándo se hunde?
    \end{itemize}
    \item \textbf{Mapa conceptual.} Construye un esquema de la unidad con estos conceptos: \emph{presión, presión hidrostática y sus factores, principio fundamental de la hidrostática, vasos comunicantes, principio de Pascal y sus aplicaciones, presión atmosférica y altitud, borrascas y anticiclones, tiempo estable e inestable, corrientes de viento, principio de Arquímedes, peso aparente, medida de densidades, flotabilidad}.
    \item \textbf{Diccionario científico.} Define con tus palabras: \emph{pascal, fluido, presión hidrostática, vasos comunicantes, prensa hidráulica, atmósfera (unidad), barómetro, mapa meteorológico, isobara, borrasca, anticiclón, corriente de viento, superficie frontal, frente frío, frente cálido, empuje, peso aparente, centro de gravedad, centro de empuje, línea de flotación, densímetro}.
\end{enumerate}

%========================================================================================
\section{Unidad 11. Trabajo y energía mecánica}
%========================================================================================
% (sacado del final de UD11_energia.tex)

\begin{enumerate}[label=\numberballred]
    \item \textbf{Resumen.} Elabora un resumen de la unidad respondiendo a estas preguntas:
    \begin{itemize}
        \item ¿Cuándo se dice que una fuerza realiza un trabajo? ¿Qué relación hay entre trabajo, fuerza y desplazamiento?
        \item ¿En qué unidades se expresan el trabajo, la energía y la potencia?
        \item ¿Qué magnitud relaciona el trabajo realizado y el tiempo empleado en hacerlo? ¿Cómo se define el rendimiento de una máquina?
        \item ¿Qué tipo de energía está asociada al movimiento de los cuerpos? ¿Y a la posición que ocupan?
        \item ¿De qué factores dependen la energía potencial gravitatoria y la energía potencial elástica?
        \item ¿Qué dice el principio de conservación de la energía mecánica? ¿Cuándo se cumple?
        \item ¿Se puede transformar una forma de energía en otra? ¿Qué dice el principio general de conservación de la energía?
    \end{itemize}
    \item \textbf{Mapa conceptual.} Construye un esquema de la unidad con estos conceptos: \emph{trabajo de una fuerza constante, potencia, potencia real y teórica, rendimiento, energía mecánica, energía cinética, energía potencial gravitatoria y elástica, el trabajo como medida de la transferencia de energía, conservación de la energía mecánica, rozamiento, transformaciones de la energía}.
    \item \textbf{Diccionario científico.} Define con tus palabras: \emph{trabajo, julio, trabajo positivo, trabajo negativo, potencia, vatio, rendimiento, kilovatio hora, caballo de vapor, electronvoltio, energía cinética, energía potencial gravitatoria, energía potencial elástica, energía mecánica, trabajo de las fuerzas de rozamiento, conservación de la energía, degradación de la energía}.
\end{enumerate}

%========================================================================================
\section{Unidad 12. El calor: una forma de transferir energía}
%========================================================================================
% (sacado del final de UD12_calor.tex)

\begin{enumerate}[label=\numberballred]
    \item \textbf{Resumen.} Elabora un resumen de la unidad respondiendo a estas preguntas:
    \begin{itemize}
        \item ¿Qué otra forma de intercambiar energía existe, además del trabajo?
        \item ¿Qué relación hay entre la energía cinética de las partículas y la temperatura?
        \item ¿De qué factores depende el calor transferido a un cuerpo cuando cambia su temperatura?
        \item ¿Qué es un calorímetro y cómo se utiliza? ¿Cómo se mide la capacidad calorífica específica de un cuerpo?
        \item ¿Qué calor se transfiere en un cambio de estado?
        \item ¿De qué depende la dilatación lineal de un sólido?
        \item ¿Qué equivalencia hay entre energía mecánica y energía térmica? ¿Se puede aprovechar toda la energía térmica como energía mecánica?
        \item ¿En qué se diferencian una máquina de combustión externa y una de combustión interna? ¿Cómo se calcula el rendimiento de una máquina térmica?
    \end{itemize}
    \item \textbf{Mapa conceptual.} Construye un esquema de la unidad con estos conceptos: \emph{intercambio de energía térmica, equilibrio térmico, cantidad de calor transferida, capacidad calorífica específica, calorímetro, calor latente de cambio de estado, dilatación de los sólidos, equivalencia entre energía mecánica y térmica, conservación de la energía total, degradación de la energía, máquinas térmicas}.
    \item \textbf{Diccionario científico.} Define con tus palabras: \emph{energía térmica, temperatura, kelvin, calor, julio, caloría, equilibrio térmico, capacidad calorífica específica, calorímetro, calor latente de fusión y de vaporización, gráfica de calentamiento, coeficiente de dilatación lineal, equivalente mecánico del calor, degradación de la energía, rendimiento, máquina térmica, máquina de vapor, motor de cuatro tiempos}.
\end{enumerate}

%========================================================================================
\section{Unidad 13. Luz y sonido: ondas que transfieren energía}
%========================================================================================
% (sacado del final de UD13_luz_sonido.tex)

\begin{enumerate}[label=\numberballred]
    \item \textbf{Resumen.} Elabora un resumen de la unidad respondiendo a estas preguntas:
    \begin{itemize}
        \item ¿Qué se propaga en un movimiento ondulatorio? ¿Cómo se clasifican las ondas?
        \item ¿Qué características tienen las ondas en un muelle, en una cuerda y en la superficie del agua?
        \item ¿Qué relación hay entre un frente de onda y un rayo?
        \item ¿Qué magnitudes caracterizan a una onda y qué relación hay entre la velocidad de propagación, la longitud de onda y el período o la frecuencia?
        \item ¿En qué consisten la reflexión y la refracción?
        \item ¿Qué características tienen las ondas sonoras? ¿Y las luminosas?
        \item ¿Qué es el nivel de intensidad sonora y en qué se mide? ¿Cuáles son las cualidades del sonido?
        \item ¿Qué es el índice de refracción? ¿Qué es el espectro electromagnético y cuáles son sus aplicaciones?
    \end{itemize}
    \item \textbf{Mapa conceptual.} Construye un esquema de la unidad con estos conceptos: \emph{ondas en muelles, cuerdas y en la superficie del agua; ondas sonoras y luminosas; criterios de clasificación; transmisión de energía sin transporte de materia; magnitudes de las ondas y relación entre ellas; velocidad de propagación; leyes de la reflexión y la refracción; intensidad física y nivel de intensidad sonora; índice de refracción; espectro electromagnético}.
    \item \textbf{Diccionario científico.} Define con tus palabras: \emph{onda (mecánica, electromagnética, longitudinal y transversal), frente de onda, rayo, amplitud, longitud de onda, período, frecuencia, hercio, intensidad, eco, decibelio, tono, timbre, umbral de audición, umbral del dolor, infrasonido, ultrasonido, imagen virtual, velocidad de la luz, índice de refracción, prisma óptico, dispersión, ondas de radio, microondas, infrarrojos, ultravioleta, rayos X, rayos gamma}.
\end{enumerate}
